% Generated draft: requires mathematical and source review.
\documentclass[11pt]{article}
\usepackage[margin=25mm]{geometry}
\usepackage{fontspec}
\setmainfont{Latin Modern Roman}
\usepackage{amsmath,amssymb,mathtools,unicode-math}
\setmathfont{Latin Modern Math}
\usepackage{enumitem,needspace,xurl,hyperref}
\setlength{\parindent}{0pt}
\setlength{\parskip}{5pt}
\setlength{\emergencystretch}{3em}
\newcommand{\N}{\mathbb{N}}
\newcommand{\Z}{\mathbb{Z}}
\newcommand{\Q}{\mathbb{Q}}
\newcommand{\R}{\mathbb{R}}
\newcommand{\C}{\mathbb{C}}
\newcommand{\F}{\mathbb{F}}
\newcommand{\A}{\mathbb{A}}
\newcommand{\PP}{\mathbb{P}}
\newcommand{\E}{\mathbb{E}}
\newcommand{\eps}{\varepsilon}
\newcommand{\dd}{\,\mathrm{d}}
\begin{document}
\section*{UnsolvedMath ID: 30002386}
\textbf{Source problem code:} catalogue entry
\subsection{Definitions and mathematical statement}
Let \(d\) be a positive integer and let \(\Z^d\) be the integer lattice. The source does not specify the intended range of dimensions. On a probability space \((\Omega,\mathcal F,\mathbb P)\), let \(W(u)\) be nonnegative random weights indexed either by vertices or by edges, according to a convention that remains unspecified. No independence, identical distribution, stationarity, nondegeneracy, or moment assumptions are supplied.

For fixed choices of admissible paths and weight-counting convention, let \(\mathcal P(x,y)\) denote the finite paths from \(x\) to \(y\), and let \(I(\gamma)\) denote the list of weighted vertices or edges traversed by \(\gamma\), with multiplicity and endpoint treatment determined by that convention. Define
\[
T(\gamma;\omega)=\sum_{u\in I(\gamma)}W(u;\omega),
\qquad
\rho(x,y;\omega)=\inf_{\gamma\in\mathcal P(x,y)}T(\gamma;\omega).
\]
Here the sum is over occurrences in the list \(I(\gamma)\). The usual nearest-neighbor path convention would require successive vertices \(v_i,v_{i+1}\) to satisfy \(\|v_{i+1}-v_i\|_1=1\); this convention is not explicitly stated in the supplied record. For vertex weights, the treatment of the initial and terminal vertices is also unspecified. Although the source calls \(\rho\) a metric, the supplied assumptions do not ensure the metric axioms; in particular, zero weights may prevent separation of distinct points.

\textbf{Source problem, with unresolved qualifications.} Determine the asymptotic order of the fluctuations of \(\rho(x,y)\) on \(\Z^d\) as the separation of the endpoints grows. The record does not specify an endpoint regime, a centering, or a quantitative meaning of fluctuation order. In particular, it does not select variance, standard deviation, quantile spread, a distributional normalization, or another notion, and it gives no numerical conjectured exponent. Thus the displayed passage-time construction records the supplied target, but a uniquely quantified asymptotic problem requires additional hypotheses and conventions. No universal fluctuation law over all nonnegative random weight fields is asserted here.
\subsection{Short English statement}
Determine the scale of random fluctuations of the minimum total passage weight between distant points of \(\Z^d\). The source leaves the dimension range, weight law, vertex or edge convention, endpoint regime, centering, and meaning of fluctuation scale unresolved. This is a statement-only draft requiring human mathematical review; the supplied open label is not an independently verified status assessment.
\subsection{Sources}
\begin{itemize}
\item Supplied local UnsolvedMath registry record, permanent ID \(30002386\), title \emph{Fluctuation Order in First-Passage Percolation}, source problem code \texttt{catalogue entry}, in \emph{Oberwolfach Reports Open Problems}. The record gives proposed year 2013 and names no proposer. Its \texttt{corrected\_verified} statement label and normalization note are supplied metadata, not independent mathematical verification.
\item Supplied report citation: \emph{Mini-Workshop: Direct and Inverse Spectral Theory of Almost Periodic Operators (2014)}, Oberwolfach Reports, pp. 2083--2117, DOI 10.4171/owr/2013/36. \url{https://doi.org/10.4171/owr/2013/36}. The literature-source list labels the same link \emph{Oberwolfach Report 36/2013}. The supplied background attributes the collection to workshop open-problem discussions. The report's exact connection to this percolation problem, and the differing year indications, require verification.
\item Supplied literature source: \emph{Lower bounds for fluctuations in first-passage percolation}. \url{https://arxiv.org/abs/1810.08270}.
\item Supplied literature source: \emph{Lower bounds on non-random fluctuations in planar first-passage percolation}. \url{https://arxiv.org/abs/2511.06166}.
\item Literature assessment dated exactly \texttt{2026-08-21}: the supplied curation labels the problem open and states that passage-time fluctuation order remains undetermined in general, with later lower-bound improvements not identifying the conjectured fluctuation exponent. This is supplied dated curation derived from the OWR collection, not an independently verified literature assessment. The source explicitly requests independent verification of time-sensitive status and cited claims.
\end{itemize}
\end{document}
